\documentclass[10pt]{article}
\usepackage[T1]{fontenc}
\usepackage[utf8]{inputenc}
\usepackage[french]{babel}
\usepackage{amsmath,amssymb,mathrsfs,mathpazo}
\usepackage{geometry}
\geometry{a4paper,margin=2cm}
\usepackage{graphicx,array,tabularx,tabularray,colortbl,multirow,multicol,booktabs,enumitem,pifont,ulem,fancybox,diagbox,cancel}
\usepackage{pstricks,pstricks-add,pst-plot,pst-tree,pst-node,pst-text,pst-func,pst-math,pst-bspline,pst-3dplot,pst-eucl,pst-solides3d}
\let\pscalculate\undefined
\usepackage{pst-bezier}
\usepackage{tikz,pgfplots,tkz-tab}
\usetikzlibrary{arrows,arrows.meta,patterns,calc,positioning,trees,decorations.markings,shapes.geometric,shapes.misc,shapes.arrows,intersections,decorations.pathmorphing,decorations.pathreplacing,backgrounds,shadows,babel}
\pgfplotsset{compat=1.18}
\usepgfplotslibrary{fillbetween}
\usepackage[np]{numprint}
\usepackage{listings,makecell,pas-tableur,twoopt}
\let\barre\undefined
\usepackage[pstricks]{tablvar}
\definecolor{darkgreen}{RGB}{0,100,0}
\providecommand{\red}{\color{red}}
\providecommand{\blue}{\color{blue}}
\providecommand{\green}{\color{green}}
\providecommand{\black}{\color{black}}
\providecommand{\white}{\color{white}}
\providecommand{\gray}{\color{gray}}
\providecommand{\cyan}{\color{cyan}}
\providecommand{\magenta}{\color{magenta}}
\providecommand{\yellow}{\color{yellow}}
\providecommand{\violet}{\color{violet}}
\providecommand{\orange}{\color{orange}}
\providecommand{\euro}{\texteuro}
\providecommand{\EUR}{\texteuro}
\providecommand{\ieme}{\textsuperscript{e}}
\providecommand{\degres}{\ensuremath{^\circ}}
\providecommand{\degree}{\ensuremath{^\circ}}
\providecommand{\vv}[1]{\overrightarrow{#1}}
\providecommand{\Coord}[1]{#1}
\providecommand{\decosix}{}
\providecommand{\starredbullet}{\ensuremath{\bullet}}
\pagestyle{empty}
\setlength{\parindent}{0pt}
\usepackage{adjustbox}\DeclareUnicodeCharacter{2236}{:}\DeclareUnicodeCharacter{2190}{\ensuremath{\leftarrow}}
\providecommand{\R}{\mathbb{R}}
\providecommand{\N}{\mathbb{N}}
\providecommand{\Z}{\mathbb{Z}}
\providecommand{\Q}{\mathbb{Q}}
\providecommand{\D}{\mathbb{D}}
\providecommand{\C}{\mathbb{C}}
\providecommand{\e}{{\mathrm{e}}}
\providecommand{\vect}[1]{\overrightarrow{#1}}
\providecommand{\vectt}[1]{\overrightarrow{\mathrm{#1}}}
\providecommand{\vvt}[1]{\overrightarrow{\mathrm{#1}}}
\providecommand{\np}[2][]{{#2~#1}}
\providecommand{\vv}[1]{{\overrightarrow{#1}}}
\renewcommand{\d}{{\mathrm{d}}}
\providecommand{\degres}{{^\circ}}
\providecommand{\degre}{{^\circ}}
\providecommand{\blue}{\color{blue}}
\providecommand{\red}{\color{red}}
\providecommand{\enskip}{\quad}
\providecommand{\llbracket}{\lbrack\!\lbrack}
\providecommand{\rrbracket}{\rbrack\!\rbrack}
\providecommand{\euro}{{\text{€}}}
\providecommand{\no}{{\text{n°}}}
\providecommand{\interval}[2]{{[#1;#2]}}
\providecommand{\intervFO}[2]{{[#1;#2[}}
\providecommand{\dotfill}{\ldots}

\providecommand{\ds}{\displaystyle}
\providecommand{\textding}[1]{#1}
\providecommand{\Oij}{\(\mathrm{O};\vec{\imath},\vec{\jmath}\)}
\providecommand{\Oijk}{\(\mathrm{O};\vec{\imath},\vec{\jmath},\vec{k}\)}
\providecommand{\Ouv}{\(\mathrm{O};\vec{u},\vec{v}\)}
\providecommand{\intervFF}[2]{\ensuremath{[#1;#2]}}
\renewcommand{\euro}{\ifmmode\text{\texteuro}\else\texteuro\fi}
\DeclareUnicodeCharacter{00B0}{\ifmmode^{\circ}\else\textdegree\fi}
\DeclareUnicodeCharacter{20AC}{\ifmmode\text{\texteuro}\else\texteuro\fi}
% Énoncés : Éducation nationale. Transcription et corrigés : APMEP, Denis Vergès et contributeurs. Découpage et classement : CoopMaths, Sébastien Lozano, Éric Elter et contributeurs. Adaptation de la présentation : Maths Explorer.
% CC BY-SA : https://coopmaths.fr/www/about/
% https://coopmaths.fr/alea/static/e3c/2020/tex/e3c_2020_00_sujet13_1_cor.tex
\begin{document}
\noindent\textit{APMEP / CoopMaths. CC BY-SA. Présentation adaptée par Maths Explorer.}\par\medskip
\section*{Exercice 1 (5 points)}
	
	\subsection*{Question 1}
La parabole a sa concavité vers le bas donc \(a < 0\) et l’équation \(g(x) = 0\) a deux solutions distinctes donc \(\Delta > 0\). \\
La réponse correcte est \textbf{c.}
	
	\subsection*{Question 2}
	
La dérivée est négative sauf sur l’intervalle \([1, 5]\) ; la fonction est donc décroissante puis croissante sur \([1, 5]\) et ensuite décroissante.\\
 La réponse correcte est \textbf{c.}
	
	\subsection*{Question 3}
	
La tangente au point d’abscisse 3 a pour équation :
\[ y - f(3) = f'(3)(x - 3). \]
Avec \(f(3) = 7\) et \(f'(3) = g(3) = 4\), une équation de la tangente est :
\[ y - 7 = 4(x - 3) \quad \text{soit} \quad y = 4x - 5. \]
La réponse correcte est \textbf{d.}
	
	\subsection*{Question 4}
	
Le vecteur \(\overrightarrow{AB} \left( \begin{array}{c} -2 \\ 3 \end{array} \right)\) est normal à tout vecteur directeur de la droite perpendiculaire. Une équation de cette droite est donc :
\[-2x + 3y + c = 0 \]
et comme le couple \((1, -3)\) vérifie cette équation, on a \(-2 - 9 + c = 0\) soit \(c = 11\).\\
 On a finalement \(-2x + 3y + 11 = 0\).\\
  La réponse correcte est \textbf{a.}
	
\subsection*{Question 5}
Par définition du produit scalaire :
	\[ \overrightarrow{BA} \cdot \overrightarrow{BC} = BA \times BC \times \cos(\widehat{BA, BC}), \]
	ce qui donne avec \(\overrightarrow{BA} \left( \begin{array}{c} 2 \\ -3 \end{array} \right)\) et \(\overrightarrow{BC} \left( \begin{array}{c} -2 \\ -5 \end{array} \right)\),
	\[ BA^2 = 4 + 9 = 13 \quad \text{et} \quad BC^2 = 4 + 25 = 29, \]
donc \(BA = \sqrt{13}\) et \(BC = \sqrt{29}\).
	
\[ \cos(\widehat{BA, BC}) = \dfrac{\overrightarrow{BA} \cdot \overrightarrow{BC}}{BA \times BC} = \dfrac{-4 + 15}{\sqrt{13} \times \sqrt{29}} = \dfrac{11}{\sqrt{13 \times 29}}. \]
La calculatrice donne \(\widehat{BA, BC} \approx 55,49^\circ\), soit au degré près $55^\circ$.\\
 La réponse correcte est \textbf{c.}

\end{document}
